\documentclass[11pt]{article}
\usepackage{mathblatt}
% gesetzt von werkzeuge/setzer.py; Lösungen nach bau/bauregeln.md „Lösungen“.
\newcommand{\kennung}{WFW}
\begin{document}
\pfheftstil
\fancyfoot[C]{{\footnotesize\color{mbgrau}\kennung\hfill\thepage}}
\pfheftkopf{Grundwert berechnen}{Klasse 7 \textperiodcentered{} Lösungen}

\begin{pfloesung}{}
\lz{1b)}{$28\,€$}{$4 \cdot 7$; $25\,\%$ passen $4$-mal in $100\,\%$}{}
\lz{c)}{$45$ m}{$5 \cdot 9$; $20\,\%$ passen $5$-mal in $100\,\%$}{}
\end{pfloesung}
\begin{pfloesung}{}
\lz{2b)}{$80 \cdot 4$}{Ganzes gesucht: der Teil passt mehrmals hinein, das Ganze ist größer; $25\,\%$ passen $4$-mal in $100\,\%$}{}
\lz{c)}{$35 \cdot 10$}{Ganzes gesucht: der Teil passt mehrmals hinein, das Ganze ist größer}{}
\lz{d)}{$0{,}1 \cdot 35$}{Teil gesucht: das Ergebnis ist kleiner als das Ganze}{}
\lz{e)}{$9 \cdot 5$}{alter Preis ist das Ganze; nicht $20\,\%$ von $9\,€$}{}
\lz{f)}{$0{,}2 \cdot 60$}{Teil gesucht: das Ergebnis ist kleiner als das Ganze}{}
\end{pfloesung}
\begin{pfloesung}{}
\lz{3b)}{$900$ kg}{$1\,\% = 27 : 3 = 9$; $100\,\% = 100 \cdot 9 = 900$}{}
\lz{c)}{$700$ m}{$1\,\% = 56 : 8 = 7$; $100\,\% = 100 \cdot 7 = 700$}{}
\lz{d)}{$600\,€$}{$1\,\% = 90 : 15 = 6$; $100\,\% = 100 \cdot 6 = 600$}{}
\lz{e)}{$30$ l}{$1\,\% = 1{,}8 : 6 = 0{,}3$; $100\,\% = 100 \cdot 0{,}3 = 30$}{}
\lz{f)}{$45$ kg}{$1\,\% = 18 : 40 = 0{,}45$; $100\,\% = 100 \cdot 0{,}45 = 45$}{}
\end{pfloesung}
\begin{pfloesung}{}
\lz{4b)}{$75$ m}{$48 : 64 \cdot 100 = 75$ m}{}
\lz{c)}{$70\,€$}{$24{,}5 : 35 \cdot 100 = 70$\,€}{}
\lz{d)}{$120\,€$}{$9{,}6 : 8 \cdot 100 = 120$\,€}{}
\lz{e)}{$180$ kg}{$81 : 45 \cdot 100 = 180$ kg}{}
\lz{f)}{$120$ m}{$3 : 2{,}5 \cdot 100 = 120$ m}{}
\end{pfloesung}
\begin{pfloesung}{}
\lz{5a)}{ja}{$1\,\% = 7 : 35 = 0{,}2$ h; $100\,\% = 20$ h $> 16$ h}{}
\lz{b)}{nein}{$1\,\% = 6 : 40 = 0{,}15$ h; $100\,\% = 15$ h $< 16$ h}{}
\end{pfloesung}
\begin{pfloesung}{}
\lz{6b)}{Prozentwert: $12$ Kinder}{Prozentwert: $0{,}3 \cdot 40 = 12$ Kinder}{}
\lz{c)}{Grundwert: $30$ Kinder}{Grundwert: $1\,\% = 18 : 60 = 0{,}3$; $100\,\% = 30$ Kinder}{}
\lz{d)}{Prozentsatz: $35\,\%$}{Prozentsatz: $84 : 240 = 0{,}35 = 35\,\%$}{}
\lz{e)}{Grundwert: $150$ Gäste}{Grundwert: $1\,\% = 135 : 90 = 1{,}5$; $100\,\% = 150$ Gäste}{}
\lz{f)}{Prozentwert: $72\,€$}{Prozentwert: $0{,}15 \cdot 480 = 72\,€$}{}
\end{pfloesung}
\begin{pfloesung}{}
\lz{7a)}{$10\,€$; $30\,€$}{Rabatt $= 0{,}25 \cdot 40 = 10\,€$; neu $= 40 - 10 = 30\,€$}{}
\lz{b)}{$20\,\%$; $72\,€$}{$18 : 90 = 0{,}2 = 20\,\%$; neu $= 90 - 18 = 72\,€$}{}
\lz{c)}{$70\,€$; $56\,€$}{$1\,\% = 14 : 20 = 0{,}70\,€$; alter Preis $= 70\,€$; neu $= 70 - 14 = 56\,€$}{}
\lz{d)}{$60\,€$; $51\,€$}{$1\,\% = 9 : 15 = 0{,}60\,€$; alter Preis $= 60\,€$; neu $= 60 - 9 = 51\,€$}{}
\end{pfloesung}
\end{document}
